% try: cite another entry of refs.bib, say \cite{lamport94}, or add a \section: numbers and references follow in one build
\documentclass[10pt,twocolumn]{article}
\usepackage[margin=0.75in,top=0.8in]{geometry}
\usepackage{amsmath,amsthm}
\usepackage{booktabs}
\usepackage{algorithm}
\usepackage{algpseudocode}

\newtheorem{theorem}{Theorem}
\newtheorem{lemma}[theorem]{Lemma}
\theoremstyle{definition}
\newtheorem{definition}[theorem]{Definition}

\title{\bfseries Keystroke-Speed Typesetting\\[2pt] \large with Exact Output}
\author{Ada Writer \and Ben Reader\\[2pt] \small Department of Typography, Example University}
\date{}

\begin{document}
\maketitle

\begin{abstract}
A writer changes one word, and a typesetter reads the whole document again.
We describe how an engine can instead keep what it built, find what an edit
can reach, and typeset only that, while every page it shows is the page a
full run would make. Section~\ref{sec:method} gives the method,
Theorem~\ref{thm:exact} its guarantee, and Table~\ref{tab:cost} where the
time goes.
\end{abstract}

\section{Introduction}
\label{sec:intro}
TeX~\cite{knuth84} and \LaTeX~\cite{lamport94} treat a document as a program
run from its first line to its last. Line breaking~\cite{knuthplass81} and
page building look only backwards, so most of a run repeats the last one
exactly. Incremental computation~\cite{acar05} suggests reusing it.

\section{Method}
\label{sec:method}
\begin{definition}
A \emph{cut} is a point in the source after which the engine's state is
determined by the state before it and the text after it.
\end{definition}
Every paragraph end is a cut. An edit between two cuts invalidates the work
from the first cut until the state agrees with the previous run again.

\begin{algorithm}[h]
\caption{Rebuild after an edit}
\label{alg:rebuild}
\begin{algorithmic}[1]
\Require cuts $c_1<\dots<c_n$, an edit at position $p$
\State $i \gets \max\{\, j : c_j \le p \,\}$
\Repeat
  \State typeset from $c_i$ to $c_{i+1}$
  \State $i \gets i + 1$
\Until{the state at $c_i$ equals the last run's}
\State reuse every page after $c_i$
\end{algorithmic}
\end{algorithm}

\begin{lemma}
\label{lem:agree}
If the state at a cut agrees with the last run's, so does everything after it.
\end{lemma}
\begin{proof}
TeX is deterministic: equal states and equal input give equal output.
\end{proof}

\begin{theorem}[Exactness]
\label{thm:exact}
The output of Algorithm~\ref{alg:rebuild} is identical, byte for byte, to a
full run on the edited source.
\end{theorem}
\begin{proof}
Up to $c_i$ the source is unchanged; from there it is typeset again until a
cut agrees, and Lemma~\ref{lem:agree} gives the rest.
\end{proof}

\section{Where the time goes}
\label{sec:cost}
\begin{table}[h]
\centering
\caption{Work per keystroke, illustrative.}
\label{tab:cost}
\begin{tabular}{@{}lrr@{}}
\toprule
Step & Full run & After an edit \\
\midrule
Read the source   & 100\% & one paragraph \\
Break into lines  & 100\% & one paragraph \\
Build the pages   & 100\% & one page \\
Write the PDF     & 100\% & one page \\
\bottomrule
\end{tabular}
\end{table}
The cost of an edit is set by what it reaches, not by the length of the
document: a word in chapter one of a thesis costs as much as a word in a
letter. The equation of the cost is simply
\begin{equation}
  T_{\text{edit}} = \sum_{k=i}^{j} t(c_k, c_{k+1}), \qquad j - i \ll n.
\end{equation}

\section{Conclusion}
Keep what is built, rebuild what an edit reaches, and prove that the result is
the full run's. The writer sees the page follow every keystroke.

\bibliographystyle{plain}
\bibliography{refs}

\end{document}
